% !TEX program = pdflatex
% ------------------------------------------------------------
% MATH 431 Homework 2
%
% Solve the following problems from the textbook:
%
% 1.12, 1.14, 1.16, 1.18, 1.32, 1.34, 1.36, 1.40
%
% Please upload your solutions here by the deadline as a single PDF file. (Have a look at the posted tips for scanning your work.) Please make sure that your submission is readable.
%
% Here is a list of additional practice exercises. You do not need to submit these problems!!
%
% 1.9, 1.11, 1.13, 1.15, 1.17, 1.19, 1.35, 1.37, 1.38, 1.39, 1.41, 1.45, 1.46, 1.48
% ------------------------------------------------------------

\documentclass[12pt]{article}

% =========================
% 0) Metadata (EDIT ME)
% =========================
\newcommand{\CourseCode}{MATH 431}
\newcommand{\CourseTitle}{Introduction to Probability}
\newcommand{\CourseSection}{LEC 002} % change to 001 or 003 if needed
\newcommand{\Semester}{Spring 2026}
\newcommand{\Instructor}{Sarah Tammen} % LEC 001: Ander Aguirre | LEC 002: Sarah Tammen | LEC 003: Nicholas Derr
\newcommand{\MeetingInfo}{MWF 9:55--10:45, 6102 Sewell Social Sciences} % LEC 001 default; update if in 002/003
\newcommand{\HomeworkNumber}{2}
\newcommand{\YourName}{Alejandro De La Torre}
\newcommand{\DueDate}{\today} % or e.g. February 2, 2026

% =========================
% 1) Core math tools
% =========================
\usepackage{amsmath, amssymb, amsthm}

% =========================
% 2) Lists and spacing
% =========================
\usepackage{enumitem}
\setlist[enumerate,1]{label=\textbf{(\alph*)}, leftmargin=2em, itemsep=0.25em}
\setlist[itemize]{leftmargin=1.5em, itemsep=0.25em}
\setlength{\parskip}{0.35em}
\setlength{\parindent}{0pt}

% =========================
% 3) Page geometry + header/footer
% =========================
\usepackage{geometry}
\geometry{margin=1in}

\usepackage{fancyhdr}
\pagestyle{fancy}
\fancyhf{}
\lhead{\CourseCode\ --- \CourseSection, \Semester}
\rhead{HW \HomeworkNumber}
\cfoot{\thepage}
\renewcommand{\headrulewidth}{0.4pt}
\setlength{\headheight}{15pt}
\addtolength{\topmargin}{-2pt}

% =========================
% 4) Links (subtle)
% =========================
\usepackage[hidelinks]{hyperref}
\setcounter{tocdepth}{1}

% =========================
% 5) Figures (optional)
% =========================
\usepackage{graphicx}
\graphicspath{{images/}}

% =========================
% 6) Lightweight boxes (optional)
% =========================
\usepackage{mdframed}
\newmdenv[
  skipabove=0.6\baselineskip,
  skipbelow=0.6\baselineskip,
  linewidth=0.6pt,
  roundcorner=4pt,
  innerleftmargin=10pt,
  innerrightmargin=10pt,
  innertopmargin=8pt,
  innerbottommargin=8pt
]{boxnote}

% =========================
% 7) Probability / notation macros
% =========================
\newcommand{\R}{\mathbb{R}}
\newcommand{\N}{\mathbb{N}}
\newcommand{\Z}{\mathbb{Z}}

\newcommand{\OmegaSpace}{\Omega}
\newcommand{\FField}{\mathcal{F}}
\newcommand{\PMeasure}{\mathbb{P}} % Probability measure in case it is relevant or want to distinguish it from P(A)
\newcommand{\E}{\mathbb{E}}

\newcommand{\given}{\,\middle|\,}
\newcommand{\ind}{\mathbf{1}}

% Common “defined as” symbol
\newcommand{\defeq}{\vcentcolon=}

% =========================
% 8) Problem command (TOC + PDF bookmarks)
% Usage:
%   \problem{<label>}
%   \problem[Short title]{<label>}
% Example:
%   \problem{1.4}
%   \problem[Sample space]{1.4}
% =========================
\makeatletter
\newcommand{\problem}[2][]{%
  \def\prob@title{Problem #2\if\relax\detokenize{#1}\relax\else\ (\textit{#1})\fi}%
  \phantomsection%
  \pdfbookmark[1]{\prob@title}{prob#2}%
  \addcontentsline{toc}{section}{\prob@title}%
  \section*{\prob@title}%
  \vspace{-0.25em}\hrule\vspace{0.8em}%
}
\makeatother

% =========================
% 9) Solution environment
% =========================
\newenvironment{solution}{%
  \noindent\textbf{Solution.}\ \ignorespaces
}{\par}

% Optional scratch / planning block
\newenvironment{scratch}{%
  \begin{boxnote}\textbf{Scratch / Setup.}\ \small
}{\end{boxnote}}

% Final answer command: expects math only (no $$ or \[ \])
\newcommand{\finalanswer}[1]{%
  \par\medskip
  \noindent\textbf{Final:}\ \fbox{$#1$}\par
} % AGAIN: MATH ONLY INSIDE THE BOX; if word answer, use \text{} inside or use \fbox{TEXT} or \framebox

% =========================
% 10) Title block
% =========================
\title{Homework \HomeworkNumber}
\author{\YourName \\
\CourseCode\ --- \CourseTitle \\
\CourseSection \\
Instructor: \Instructor}
\date{Due: February 6, 2026}

\begin{document}
\maketitle

% \section*{Collaboration / Resources}
% (Optional) Write “I worked alone” or list collaborators/resources.

\tableofcontents
\bigskip
\hrule
\bigskip
\newpage
%==========================================================
% Problems
%==========================================================

\problem{1.12}
We roll a fair die repeatedly until we see the number four appear and then we stop.
\begin{enumerate}
\item What is the probability that we need at most 3 rolls?
\item What is the probability that we needed an even number of die rolls?
\end{enumerate}

\begin{solution}
\begin{enumerate}
\item Let $N$ be the roll on which the first 4 appears. Then $N$ is geometric with $p=1/6$.
\[
P(N \le 3)=1-P(\text{no 4 in first 3 rolls})=1-\left(\frac{5}{6}\right)^3.
\]
\finalanswer{1-\left(\frac{5}{6}\right)^3}
\item
\[
P(N \text{ is even})=\sum_{k=1}^\infty P(N=2k)=\sum_{k=1}^\infty \left(\frac{5}{6}\right)^{2k-1}\frac{1}{6}.
\]
This is a geometric series with ratio $(5/6)^2$, so
\[
P(N \text{ is even})=\frac{\frac{1}{6}\cdot \frac{5}{6}}{1-\left(\frac{5}{6}\right)^2}=\frac{5}{11}.
\]
\finalanswer{\frac{5}{11}}
\end{enumerate}
\end{solution}
\newpage

\problem{1.14}
Assume that $P(A) = 0.4$ and $P(B) = 0.7$. Making no further assumptions on $A$ and $B$, show that $P(AB)$ satisfies $0.1 \le P(AB) \le 0.4$.

\begin{solution}
By inclusion--exclusion,
\[
P(A\cup B)=P(A)+P(B)-P(AB)\le 1,
\]
so $P(AB)\ge P(A)+P(B)-1=0.4+0.7-1=0.1$. Also,
\[
P(AB)\le \min\{P(A),P(B)\}=\min\{0.4,0.7\}=0.4.
\]
\finalanswer{0.1\le P(AB)\le 0.4}
\end{solution}
\newpage

\problem{1.16}
We flip a fair coin five times. For every heads you pay me \$1 and for every tails I pay you \$1. Let $X$ denote my net winnings at the end of five flips. Find the possible values and the probability mass function of $X$.

\begin{solution}
Let $H$ be the number of heads in 5 flips. Then $H\sim \text{Binomial}(5,1/2)$ and
\[
X=(\#\text{heads})-(\#\text{tails})=H-(5-H)=2H-5.
\]
Thus the possible values are
\[
X\in\{-5,-3,-1,1,3,5\},
\]
with probability mass function
\[
P(X=2h-5)=P(H=h)=\binom{5}{h}2^{-5},\qquad h=0,1,2,3,4,5.
\]
\finalanswer{P(X=2h-5)=\binom{5}{h}/2^{5},\ h=0,\dots,5}
\end{solution}
\newpage

\problem{1.18}
The statement
\begin{center}
SOME DOGS ARE BROWN
\end{center}
has 16 letters. Choose one of the 16 letters uniformly at random. Let $X$ denote the length of the \emph{word} containing the chosen letter. Determine the possible values and probability mass function of $X$.

\begin{solution}
The words and their lengths are: SOME (4), DOGS (4), ARE (3), BROWN (5). Since 16 letters are equally likely,
\[
P(X=3)=\frac{3}{16},\quad P(X=4)=\frac{8}{16}=\frac12,\quad P(X=5)=\frac{5}{16}.
\]
\finalanswer{X\in\{3,4,5\},\ P(X=3)=\frac{3}{16},\ P(X=4)=\frac12,\ P(X=5)=\frac{5}{16}}
\end{solution}
\newpage

\problem{1.32}
You are dealt five cards from a standard deck of 52. Find the probability of being dealt a \emph{full house}. (This means that you have three cards of one face value and two cards of a different face value. An example would be three queens and two fours. See Exercise C.10 in Appendix C for a description of all the poker hands.)
\begin{solution}
Count full houses and divide by total hands. Choose the rank of the triple, then its suits, then the rank and suits of the pair:
\[
N_{\text{fh}}=13\binom{4}{3}\cdot 12\binom{4}{2}=13\cdot 4\cdot 12\cdot 6=3744.
\]
Total 5-card hands are $\binom{52}{5}$, so
\[
P(\text{full house})=\frac{3744}{\binom{52}{5}}=\frac{6}{4165}.
\]
\finalanswer{\frac{6}{4165}}
\end{solution}
\newpage

\problem{1.34}
Pick a uniformly chosen random point inside a unit square (a square of sidelength 1) and draw a circle of radius $1/3$ around the point. Find the probability that the circle lies entirely inside the square.
\begin{solution}
For the circle of radius $1/3$ to lie inside the unit square, the center must be at least $1/3$ from each side. Thus the center must lie in
\[
\left[\frac13,\frac23\right]\times\left[\frac13,\frac23\right],
\]
a square of side length $1-2\cdot\frac13=\frac13$, area $1/9$. By uniformity,
\[
P(\text{circle inside})=\frac{1}{9}.
\]
\finalanswer{\frac{1}{9}}
\end{solution}
\newpage

\problem{1.36}
\begin{enumerate}
\item Let $(X, Y)$ denote a uniformly chosen random point inside the unit square
\[
[0,1]^2 = [0,1] \times [0,1] = \{(x,y): 0 \le x,y \le 1\}.
\]
Let $0 \le a < b \le 1$. Find the probability $P(a < X < b)$, that is, the probability that the $x$-coordinate $X$ of the chosen point lies in the interval $(a,b)$.
\item What is the probability $P(|X - Y| \le 1/4)$?
\end{enumerate}
\begin{solution}
\begin{enumerate}
\item The event $\{a<X<b\}$ corresponds to the vertical strip $(a,b)\times[0,1]$, which has area $b-a$. Since the point is uniform,
\[
P(a<X<b)=b-a.
\]
\finalanswer{b-a}
\item The condition $|X-Y|\le \frac14$ is the band between the lines $y=x\pm \frac14$ in the unit square. The excluded regions are two right triangles, each with legs of length $3/4$, so each has area $\frac12\left(\frac34\right)\left(\frac34\right)=\frac{9}{32}$. Hence
\[
P(|X-Y|\le \tfrac14)=1-2\cdot\frac{9}{32}=1-\frac{9}{16}=\frac{7}{16}.
\]
\finalanswer{\frac{7}{16}}
\end{enumerate}
\end{solution}
\newpage

\problem{1.40}
An urn contains 1 green ball, 1 red ball, 1 yellow ball and 1 white ball. I draw 4 balls with replacement. What is the probability that there is at least one color that is repeated exactly twice?

Hint. Use inclusion--exclusion with events $G = \{\text{exactly two balls are green}\}$, $R = \{\text{exactly two balls are red}\}$, etc.
\begin{solution}
There are $4^4=256$ equally likely sequences. Let $G$ be the event that exactly two balls are green, and similarly $R,Y,W$.
For $G$, choose the two green positions and fill the other two positions with any of the other 3 colors:
\[
|G|=\binom{4}{2}3^2=6\cdot 9=54.
\]
Thus $\sum |G|=4\cdot 54=216$. For a pair, say $G\cap R$, we need exactly two greens and two reds, so
\[
|G\cap R|=\binom{4}{2}=6,
\]
and there are $\binom{4}{2}=6$ such pairs, giving $\sum |G\cap R|=36$. No triple intersections are possible. By inclusion--exclusion,
\[
|\cup G|=216-36=180.
\]
Therefore
\[
P(\text{at least one color appears exactly twice})=\frac{180}{256}=\frac{45}{64}.
\]
\finalanswer{\frac{45}{64}}
\end{solution}


\end{document}
